\subsection{Psychisches Wohlbefinden}   \label{subsec_2.1.3}
Die Definition des mehrdimensionalen psychischen Wohlbefindens lässt sich nicht leicht aufstellen. Unterschiedliche Wissenschaftler im Bereich der Psychologie, Phylisophie und Sozialwissenschaften haben ihre Definitionen aufgestellt, die nicht alle deckungsgleich sind, im Kern jedoch auf einen gemeinsamen Nennen treffen \parencite{PsychWohlbefinden}. 

Wenn man sich mit dem Begriff des Wohlbefindens beschäftigt, muss man sich mit einem vieldeutigen Konzept auseinandersetzen, das unterschiedliche disziplinäre und definitorische Hintergründe besitzt. \parencite{Wohlbefinden}










In den Sozialwissenschaften und der Medizin zählt das psychische Wohlbefinden als Teilaspekt von Gesundheit allgemein und spezifischer der psychischen Gesundheit. % Buch 2001 in Bachelor-Modul Gesundheitspsychologie I-Literatur

Definition
Unterschied Wohlbefinden und Lebenszufriedenheit/ Lebensqualität (gilt als kognitiver "Trait" als Teil von subjektivem Wohlbefinden)
Wie im Rahmen dieser Studie zu verstehen (WHO-5 erhebt subjektive Perspektive)
Modell zur Aufrechterhaltung
Bio-Psycho-Soziale Modell



%\subsubsection{Vergewaltigungsmythen}  \label{2.1.3.1}
%\subsubsection{Victim Blaming}     \label{2.1.3.2}
%Kapitel ~\ref{2.1.3.2} 
%\subsubsection{Theorien zur Erklärung von Victim Blaming}     \label{2.1.3.3}
